\section*{ЗАКЛЮЧЕНИЕ}
\addcontentsline{toc}{section}{ЗАКЛЮЧЕНИЕ}

В ходе выполнения выпускной квалификационной работы была рассмотрена задача оценки кредитоспособности физических лиц в банковских системах с использованием алгоритма градиентного бустинга деревьев решений.

В аналитическом разделе был проведен анализ предметной области кредитного скоринга. Были рассмотрены основные понятия, связанные с оценкой кредитного риска, дефолтом заемщика, скоринговым баллом и процессом принятия кредитного решения. Также были проанализированы традиционные статистические методы и современные алгоритмы машинного обучения, применяемые в задачах кредитного скоринга. В результате анализа было обосновано использование алгоритма градиентного бустинга деревьев решений как метода, способного учитывать сложные нелинейные зависимости между признаками заемщика и вероятностью дефолта.
Была сформулирована цель работы и выполнена формализация задачи оценки кредитоспособности с использованием IDEF0-диаграмм. В результате были выделены основные элементы процесса оценки кредитоспособности: исходная информация о заемщике, правила и ограничения, используемые средства обработки данных, а также итоговые результаты скоринговой оценки.

В технологическом разделе описан программный комплекс для оценки кредитоспособности физических лиц. В состав программного комплекса входят модуль формирования набора данных, модуль обучения модели, модуль предсказания кредитоспособности и пользовательский интерфейс. Для реализации модулей обработки данных и машинного обучения был использован язык Python, а для разработки графического интерфейса --- язык C\# и технология Windows Forms.

В рамках разработки были определены форматы входных и выходных данных. Входные данные представляются в формате JSON и содержат финансовые, социально-демографические характеристики заемщика, сведения о кредитной истории и агрегированные показатели платежного поведения. Выходные данные включают вероятность дефолта, скоринговый балл, уровень кредитного риска, решение модели, факторы риска и рекомендации. Проведено тестирование программного обеспечения на тестовой выборке и синтетических наборах данных. В ходе тестирования проверены основные функции программного комплекса: ручной ввод данных, формирование JSON-структуры, запуск модуля предсказания из приложения, получение результата скоринга и вывод итоговой оценки в интерфейсе. Результаты тестирования подтвердили корректность работы основных компонентов программного комплекса.

В исследовательском разделе была проведена оценка эффективности разработанного метода. Для сравнения были рассмотрены логистическая регрессия, случайный лес и модель градиентного бустинга LightGBM. Экспериментальные результаты показали, что модель LightGBM обеспечивает наилучшее качество среди рассмотренных алгоритмов по метрикам ROC-AUC и F1-score. Было установлено, что в условиях дисбаланса классов метрика Accuracy не является достаточной для оценки качества модели. Использование пониженного порога классификации позволило повысить полноту выявления заемщиков с высоким риском дефолта. При этом снижение порога привело к увеличению числа ложноположительных решений, что является характерным компромиссом для задач кредитного скоринга.

Был выполнен анализ значимости признаков модели. Наибольшее влияние на результат оказали внешние скоринговые оценки, возраст заемщика, сумма кредита, размер ежемесячного платежа, стаж занятости и признаки, связанные с предыдущими заявками. Полученные результаты соответствуют предметной области кредитного скоринга и подтверждают, что модель использует содержательно значимые факторы кредитного риска.

Таким образом, поставленная цель работы была достигнута: был разработан метод оценки кредитоспособности физических лиц на основе алгоритма градиентного бустинга деревьев решений и создан программный комплекс, реализующий данный метод. Разработанное программное средство может использоваться как инструмент предварительной оценки кредитного риска заемщика, при этом итоговое решение о выдаче кредита должно приниматься с учетом дополнительной проверки и экспертной оценки специалиста.
